\documentclass[12pt,a4paper]{article}

\usepackage[margin=2.5cm]{geometry}
\usepackage{amsmath, amssymb, amsthm}
\usepackage{booktabs}
\usepackage{array}
\usepackage{tabularx}
\usepackage{enumitem}
\usepackage{xcolor}
\usepackage{titlesec}
\usepackage{parskip}
\usepackage{fancyhdr}
\usepackage{hyperref}
\usepackage{tcolorbox}
\usepackage{mathtools}
\usepackage{bm}

% Colors
\definecolor{sectioncolor}{RGB}{180,0,0}
\definecolor{boxbg}{RGB}{245,245,245}
\definecolor{boxframe}{RGB}{180,0,0}
\definecolor{notebg}{RGB}{255,250,235}
\definecolor{noteborder}{RGB}{200,150,0}

% Section formatting
\titleformat{\section}{\large\bfseries\color{sectioncolor}}{}{0em}{}[\vspace{0.2em}{\color{sectioncolor}\hrule height 1pt}]
\titleformat{\subsection}{\normalsize\bfseries}{}{0em}{}
\titleformat{\subsubsection}{\normalsize\itshape}{}{0em}{}

% Header/footer
\pagestyle{fancy}
\fancyhf{}
\rhead{\textit{Imran Sir --- ML Notes by Hanium}}
\lhead{}
\rfoot{\thepage}
\renewcommand{\headrulewidth}{0.4pt}

% Box environments
\tcbuselibrary{skins, breakable}
\newtcolorbox{definitionbox}[1][]{
  colback=boxbg, colframe=boxframe,
  fonttitle=\bfseries, title=#1,
  breakable, left=6pt, right=6pt, top=4pt, bottom=4pt
}
\newtcolorbox{examplebox}{
  colback=notebg, colframe=noteborder,
  fonttitle=\bfseries, title=Example,
  breakable, left=6pt, right=6pt, top=4pt, bottom=4pt
}
\newtcolorbox{notebox}{
  colback=notebg, colframe=noteborder,
  breakable, left=6pt, right=6pt, top=4pt, bottom=4pt
}

\newcommand{\keyword}[1]{\textbf{#1}}
\newcommand{\algo}[1]{\textbf{\textit{#1}}}

\title{\Huge\bfseries\color{sectioncolor} Machine Learning Notes\\[0.3em]
       \large\normalfont\color{black} Imran Sir \quad $\bullet$ \quad Written by Hanium}
\date{}

\begin{document}
\maketitle
\tableofcontents
\newpage

%==============================================================
\section{K-Means Clustering}
%==============================================================

\begin{definitionbox}[Definition]
K-Means is an \keyword{unsupervised iterative clustering technique}. A \keyword{cluster} is defined as a collection of data points exhibiting certain similarities.
\end{definitionbox}

\subsection{Steps}
\begin{enumerate}[leftmargin=*, label=\textbf{Step \arabic*.}]
  \item Choose the number of clusters $k$.
  \item Randomly select any $k$ data points as \textbf{cluster centers}.
  \item Calculate the distance between each data point and each cluster center.
  \item Assign each data point to the nearest cluster (nearest center).
  \item Re-compute the center of each newly formed cluster.
  \item Repeat Steps 3--5 until one of the following stopping conditions is met:
        \begin{itemize}
          \item Centers of new clusters do not change.
          \item Data points remain in the same cluster.
          \item Maximum iteration number is reached.
        \end{itemize}
\end{enumerate}

\subsection{Advantages}
\begin{itemize}
  \item Iteratively efficient: $O(nkt)$, where $n$ = points, $k$ = clusters, $t$ = iterations.
  \item Often terminates at local optimum.
\end{itemize}

\subsection{Disadvantages}
\begin{itemize}
  \item Sensitive to outliers.
  \item Requires the number of clusters $k$ to be specified in advance.
  \item Cannot handle noisy data and outliers well.
  \item Not suitable for non-convex shapes.
  \item Random centroid choosing may give poor results.
\end{itemize}

\subsection{Example}

\begin{examplebox}
Data points: $2, 4, 6, 10, 12, 14$ \quad (1-D)\par
Let $k = 2$, initial centroids: $C_1 = 4$, $C_2 = 12$.

\vspace{0.5em}
\textbf{Iteration 1 --- Assignment:}
\begin{center}
\begin{tabular}{cccc}
\toprule
Point & dist from $C_1$ & dist from $C_2$ & Cluster \\
\midrule
2 & 2 & 10 & $C_1$ \\
4 & 0 & 8  & $C_1$ \\
6 & 2 & 6  & $C_1$ \\
10 & 6 & 2 & $C_2$ \\
12 & 8 & 0 & $C_2$ \\
14 & 10 & 2 & $C_2$ \\
\bottomrule
\end{tabular}
\end{center}

Cluster 1: $\{2,4,6\}$, Cluster 2: $\{10,12,14\}$

\textbf{Update centroids:}
\[
C_1 = \frac{2+4+6}{3} = 4, \qquad C_2 = \frac{10+12+14}{3} = 12
\]
Centroids unchanged $\Rightarrow$ \textbf{Converged.}
\end{examplebox}

%==============================================================
\section{Finding the Value of \texorpdfstring{$k$}{k} --- Elbow Method}
%==============================================================

\begin{definitionbox}[WCSS --- Within Cluster Sum of Squares]
For each $k$, we compute:
\[
\text{WCSS} = \sum_{C_k} \sum_{d_i \in C_k} \bigl(\text{dist}(d_i, C_k)\bigr)^2
\]
where $C_k$ is the cluster centroid and $d_i$ are data points in each cluster.
\end{definitionbox}

Plot WCSS vs.\ $k$. The \textbf{elbow point} (where the rate of decrease slows sharply) gives the optimal $k$.

\[
\text{WCSS (example values): } 8000, 7000, 6000, 5000, \ldots
\]

\subsection{Problems with the Elbow Method} \quad \textit{[Complexity: $O(nkd)$]}

\begin{itemize}
  \item The lower the WCSS, the better the clustering --- but this is \textbf{not always true}.
  \item If $k$ is very large, WCSS will be reduced but clustering is not good at all. In the extreme case it becomes one cluster per point.
  \item \textbf{That's why the Silhouette Coefficient is preferred over the Elbow method.}
  \item Comparison: Elbow method complexity $= O(nkd)$; Silhouette $= \dfrac{n^2d}{nkd} = \dfrac{1}{k}$ times $\Rightarrow$ silhouette is better.
\end{itemize}

%==============================================================
\section{Silhouette Method}
%==============================================================

\begin{definitionbox}[Silhouette Method]
Used to find how good the clustering is, and to choose the best value of $k$ in K-Means. It checks:
\begin{itemize}
  \item How close a point is to its own cluster.
  \item How far it is from other clusters.
  \item Complexity: $O(n^2 k)$.
\end{itemize}
\end{definitionbox}

\subsection{Silhouette Coefficient Formula}

\[
\boxed{S(i) = \frac{b(i) - a(i)}{\max\{a(i),\, b(i)\}}, \qquad -1 \leq S \leq 1}
\]

where:
\begin{align*}
a(i) &= \text{average distance from point } i \text{ to other points in the \textbf{same} cluster} \\
b(i) &= \text{average distance from point } i \text{ to points in the \textbf{nearest other} cluster}
\end{align*}

\subsection{Interpretation Table}

\begin{center}
\begin{tabular}{cc}
\toprule
$S$ value & Interpretation \\
\midrule
$+1$ (near) & Good clustering \\
$0$         & Overlapping clusters \\
$-1$        & Wrong cluster assignment \\
\bottomrule
\end{tabular}
\end{center}

\subsection{Example}

\begin{examplebox}
Two clusters: $A = \{A_1(1,1),\; A_2(2,1),\; A_3(1,2),\; A_4(2,2)\}$,\quad $B = \{B_1(8,8),\; B_2(9,8),\; B_3(8,9),\; B_4(9,9)\}$

\textbf{Step 1:} Choose point $P = A_1 = (1,1)$.

\textbf{Step 2:} Calculate $a(i) = a(P) = a(1,A)$ (average distance to other 3 points in $A$):
\[
a_i = \frac{d_{12} + d_{13} + d_{14}}{3} = \frac{1 + 1 + \sqrt{2}}{3} \approx 1.14
\]

\textbf{Step 3:} Calculate $b(i)$ --- average distance to all 4 points in nearest other cluster $B$:
\[
b(i) = \frac{d_{11}+d_{12}+d_{13}+d_{14}}{4} \approx 10.62
\]

\textbf{Step 4:}
\[
S(i) = \frac{b(i) - a(i)}{\max\{a(i), b(i)\}} = \frac{10.62 - 1.14}{10.62} \approx 0.89 \quad \text{(very good clustering)}
\]
\end{examplebox}

%==============================================================
\section{K-Means++ Clustering}
%==============================================================

\begin{definitionbox}[K-Means++]
An \textbf{improved version} of K-Means clustering. K-Means chooses centroids randomly, which may lead to poor clustering. K-Means++ chooses \textbf{better initial centroids}.

\[
P(x_i) = \frac{D(x_i)}{\sum D(x_i)}
\]
where $D(x_i)$ is the squared distance from point $x_i$ to the nearest already-chosen centroid.
\end{definitionbox}

\subsection{Why K-Means++ is Probabilistic}

\[
P(x_i) = \frac{D(x_i)}{\sum D(x_i)} \qquad \text{(probabilistic / random feature)}
\]

K-Means++ chooses the first centroid randomly, then subsequent centroids probabilistically. So different runs may produce different centroids and clusters. This helps avoid:
\begin{itemize}
  \item Bad initialization
  \item Local minima
  \item Biased centroid selection
\end{itemize}

\subsubsection{What if deterministic?}
\begin{enumerate}
  \item Can get stuck in poor solution.
  \item Centroids will be wasted on noisy data (a noisy point would become a centroid if deterministic, but probabilistically this is unlikely).
  \item The same centroid won't be chosen again because $P(\text{already chosen}) = 0$ (from the equation).
\end{enumerate}

\subsection{Example --- K-Means++ with $k=2$}

\begin{examplebox}
Points: $(1,1),(2,1),(1,2),(2,2),(8,8),(8,9),(9,8),(9,9)$, $k=2$.

\textbf{Step 1:} Choose first centroid randomly: $C_1 = (1,1)$.

\textbf{Step 2:} Compute $D(x_i) = \|x_i - C_1\|^2$:
\begin{center}
\begin{tabular}{cc}
\toprule
Point & $D(x_i)$ \\
\midrule
$(1,1)$ & 0 \\
$(2,1)$ & 1 \\
$(1,2)$ & 1 \\
$(2,2)$ & 2 \\
$(8,8)$ & 98 \\
$(8,9)$ & 113 \\
$(9,8)$ & 113 \\
$(9,9)$ & 128 \\
\bottomrule
\end{tabular}
\end{center}

$\sum D(x_i) = 456$. Points far away have higher probability.
\[
p = \frac{D(x_i)}{\sum D(x_i)}
\]
Maximum probability at $(9,9)$, so $C_2 = (9,9)$.

\textbf{Step 4:} Repeat K-Means algorithm with these two good centroids.
\end{examplebox}

\subsection{Example --- K-Means++ with $k=3$}

\begin{examplebox}
Points: $(1,1),(2,1),(1,2),(2,2),(8,8),(9,8),(8,9),(20,20),(21,20),(20,21)$, $k=3$.

\textbf{Step 1:} Randomly choose $C_1 = (1,1)$. Compute $D(x_i)$.
Last three points have highest squared distances: $(20,20)\to 722$, $(21,20)\to 761$, $(20,21)\to 761$.

Let $C_2 = (20,20)$.

\textbf{Step 2:} Find third centroid. Now $D(x) = \min(d(x,C_1),\; d(x,C_2))$:
\begin{center}
\begin{tabular}{cccc}
\toprule
Point & $d(x,C_1)$ & $d(x,C_2)$ & $D$ \\
\midrule
$(1,1)$ & 0 & 722 & 0 \\
$(2,1)$ & 1 & 685 & 1 \\
$(1,2)$ & 1 & 685 & 1 \\
$(2,2)$ & 2 & 648 & 2 \\
$(8,8)$ & 98 & 288 & 98 \\
$(8,9)$ & 113 & 265 & 113 \\
$(9,8)$ & 113 & 265 & 113 \\
$(20,20)$ & 722 & 0 & 0 \\
$(21,20)$ & 761 & 1 & 1 \\
$(20,21)$ & 761 & 1 & 1 \\
\bottomrule
\end{tabular}
\end{center}
Higher probabilities at $(8,8),(8,9),(9,8)$. Let $C_3 = (8,8)$.

\textbf{Step 3:} Run K-Means algorithm.
\end{examplebox}

\subsection{Advantages and Disadvantages of K-Means++}

\begin{minipage}[t]{0.48\textwidth}
\textbf{Advantages:}
\begin{itemize}
  \item Better initial centroids
  \item Lower WCSS
  \item Less computation
  \item Probabilistic (different solutions in different approaches)
\end{itemize}
\end{minipage}
\hfill
\begin{minipage}[t]{0.48\textwidth}
\textbf{Disadvantages:}
\begin{itemize}
  \item Chance of getting stuck in local minima
\end{itemize}
\end{minipage}

%==============================================================
\section{Genetic Algorithm (for K-Means)}
%==============================================================

\begin{definitionbox}[Genetic Algorithm]
The local-minima problem of K-Means (and K-Means++) can be solved using a Genetic Algorithm.
\end{definitionbox}

\subsection{Steps}

\begin{enumerate}[leftmargin=*, label=\textbf{\arabic*.}]
  \item \textbf{Initial Population:}
        We need a \textbf{solution set} (e.g.\ 20 solutions total):
        \[
        [l_1, l_3, l_5],\quad [l_2, l_4, l_8], \ldots
        \]
        To make the solution set, run K-Means algorithm 20\%. Others are made randomly.

  \item \textbf{Selection:}
        \begin{itemize}
          \item Define \textbf{fitness}: lower WCSS $\Rightarrow$ better solution.
          \item For betterment, first run K-Means, then WCSS.
          \item Methods: Roulette wheel, Tournament selection.
        \end{itemize}

  \item \textbf{Mutation:}
        \begin{itemize}
          \item Randomly modify centroid.
          \item We choose the point for mutation which is causing larger WCSS.
        \end{itemize}

  \item \textbf{Elitism:}
        \begin{itemize}
          \item Will store the best solutions.
        \end{itemize}
\end{enumerate}

%==============================================================
\section{Non-Convex Data Clustering}
%==============================================================

K-Means relies on Euclidean distance to assign points to the nearest centroid. However, it cannot perfectly partition non-convex data:
\begin{itemize}
  \item It splits a single continuous non-convex cluster into multiple \textbf{artificial convex regions}.
  \item It merges separate but interleaved non-convex structures into a single incorrect group.
\end{itemize}

\textbf{That's why we use DB-SCAN.}

%==============================================================
\section{DBSCAN}
%==============================================================

\begin{definitionbox}[DBSCAN --- Density-Based Spatial Clustering of Applications with Noise]
\begin{itemize}
  \item A clustering algorithm that groups points \textbf{based on density}.
  \item Works for non-convex shapes. \hfill \textit{(Solves two problems)}
  \item Detects outliers on noisy data points.
  \item \textbf{Two parameters:}
        \begin{enumerate}
          \item $\varepsilon$ (epsilon): radius around a point.
          \item \texttt{minPts}: minimum points needed to form a region ($= 2d$, where $d$ = dimensions).
        \end{enumerate}
\end{itemize}
\end{definitionbox}

\subsection{Three Types of Points}
\begin{enumerate}
  \item \textbf{Core}: has $\geq \texttt{minPts}$ within $\varepsilon$.
  \item \textbf{Border}: does not have \texttt{minPts} neighbors, but is within $\varepsilon$ of a core point.
  \item \textbf{Noise}: not close enough to any dense point.
\end{enumerate}

\subsection{Example}

\begin{examplebox}
Points: $(1,1),(1,2),(2,1),(2,2),(8,8),(8,9),(9,8),(9,9),(20,20),(25,25)$\\
$\texttt{minPts} = 2 \times 2 = 4$

\textbf{Distance matrix} (10$\times$10 shown partially):

\begin{center}
\begin{tabular}{c|cccccccccc}
 & p1 & p2 & p3 & p4 & p5 & p6 & p7 & p8 & p9 & p10 \\
\hline
p1 & 0 & 1 & 1 & 1.41 & \textbf{9.9} & 10.63 & 10.63 & 11.31 & 26.9 & 33.9 \\
p2 & 1 & 0 & 1.41 & 1 & \textbf{9.22} & 9.9 & 10 & 10.63 & 26.17 & 33.24 \\
p3 & 1 & 1.41 & 0 & 1 & \textbf{9.21} & 10 & 9.9 & 10.63 & 26.17 & 33.24 \\
p4 & 1.41 & 1 & 1 & 0 & \textbf{8.5} & 9.22 & 9.22 & 9.9 & 25.5 & 32.5 \\
p5 & 9.9 & 9.22 & 9.21 & 8.5 & 0 & 1 & 1 & 1.41 & 16.9 & \textbf{8.24} \\
p6 & 10.63 & 9.9 & 10 & \textbf{9.22} & 1 & 0 & 1.41 & 1 & 16.3 & 23.3 \\
p7 & 10.63 & 10 & 9.9 & \textbf{9.22} & 1 & 1.41 & 0 & 1 & 16.24 & 23.3 \\
p8 & 11.31 & 10.63 & 10.63 & \textbf{9.9} & 1.41 & 1 & 1 & 0 & 6.63 & 22.63 \\
p9 & 26.9 & 26.17 & 26.17 & 25.5 & 16.9 & 16.3 & 16.24 & 6.63 & 0 & 7.07 \\
p10 & 33.9 & 33.24 & 33.24 & 32.5 & 8.24 & 23.3 & 23.3 & 22.63 & 7.07 & 0 \\
\end{tabular}
\end{center}

\textbf{$\varepsilon$-determination:} For each point, sort distances; the ``elbow'' value gives $\varepsilon$.

Sorted (4th-nearest distances): $8.24, 8.5, 9.21, 9.22, 9.22, 9.22, 9.9, 9.9, 16.3, 23.3$

The elbow occurs at $\varepsilon \approx 9.22$.

\textbf{Considering $\varepsilon = 9.22$:}
\begin{itemize}
  \item p1: neighbors $\{$p2, p3, p4, p5$\}$ $\to$ \textbf{Core}
  \item p2: neighbors $\{$p1, p4, p3, p5, p6$\}$ $\to$ \textbf{Core}
  \item p3: neighbors $\{$p1, p2, p4, p5, p7$\}$ $\to$ \textbf{Core}
  \item p4: neighbors $\{$p1, p2, p3, p5, p6, p7$\}$ $\to$ \textbf{Core}
  \item p5: neighbors $\{$p1, p2, p3, p4, p6, p7, p8, p10$\}$ $\to$ \textbf{Core}
  \item p6: neighbors $\{$p2, p4, p5, p7, p8$\}$ $\to$ \textbf{Core}
  \item p7: neighbors $\{$p3, p4, p5, p6, p8$\}$ $\to$ \textbf{Core}
  \item p8: neighbors $\{$p4, p5, p6, p7, p9$\}$ $\to$ \textbf{Core}
  \item p9: neighbors $\{$p5, p6$+1, p10$\}$ $\to$ \textbf{Border}
  \item p10: neighbors $\{$p5, p9$\}$ $\to$ \textbf{Border}
\end{itemize}
\end{examplebox}

\subsection{Advantages and Disadvantages of DBSCAN}

\begin{minipage}[t]{0.48\textwidth}
\textbf{Advantages:}
\begin{itemize}
  \item Defining $\varepsilon$ and \texttt{minPts} is intuitive (neutral)
  \item Works for non-convex data clustering
\end{itemize}
\end{minipage}
\hfill
\begin{minipage}[t]{0.48\textwidth}
\textbf{Disadvantages:}
\begin{itemize}
  \item $\varepsilon$-dependency
  \item Cannot detect varying distance
  \item Non-deterministic
  \item Large $\varepsilon$ $\to$ single cluster; small $\varepsilon$ $\to$ many clusters
  \item Cannot work if various densities are present
  \item One $\varepsilon$ may not work for clusters with different densities
  \item Single $\varepsilon$ problem:
        \begin{itemize}
          \item $\varepsilon$ small: dense cluster recovered but sparse cluster may fragment
          \item $\varepsilon$ large: bridge points can connect two clusters
        \end{itemize}
\end{itemize}
\end{minipage}

\vspace{1em}
\begin{notebox}
\textbf{WCAD (Curse of Dimensionality):} Let $\eta = 0.5$. If $d = 2$: $(0.5)^2 = 0.25 = 25\%$ of data. If $d = 10$: $(0.5)^{10} < 1\%$ of data.\\
$\Rightarrow$ \textbf{The more the dimension, the more sparse the data will be.}
\end{notebox}

%==============================================================
\section{H-DBSCAN (Hierarchical DBSCAN)}
%==============================================================

\begin{definitionbox}[H-DBSCAN]
\begin{itemize}
  \item \textbf{Hierarchical DBSCAN} --- an improved version of DBSCAN.
  \item Instead of using one fixed $\varepsilon$, it tries many density levels, builds a \textbf{hierarchy of clusters}, and keeps the most stable clusters.
  \item Handles \textbf{variable density}.
\end{itemize}
\end{definitionbox}

\subsection{Algorithm}
\begin{enumerate}
  \item Calculate distance matrix.
  \item Calculate core distance.
  \item Compute mutual reachability distance (MRD).
  \item Construct MST (Prim, Kruskal).
  \item Condense the tree and identify clusters.
  \item Output final clusters.
\end{enumerate}

\subsection{H-DBSCAN Example}

\begin{examplebox}
Points: $A(1,2),\; B(2,2),\; C(2,3),\; D(7,8),\; E(8,7),\; F(25,80)$\quad $\texttt{minPts} = 2$

\textbf{Step 1: Distance Matrix}
\begin{center}
\begin{tabular}{ccccccc}
\toprule
 & A & B & C & D & E & F \\
\midrule
A & 0    & \textbf{1}    & \textbf{1.41} & 8.5  & 8.6  & 82.5  \\
B & 1    & 0    & \textbf{1}    & 7.81 & 7.81 & 82.34 \\
C & 1.41 & 1    & 0    & 7.07 & 7.21 & 80.33 \\
D & 8.5  & 7.81 & 7.07 & 0    & 1.41 & 74.19 \\
E & 8.6  & 7.81 & 7.21 & 1.41 & 0    & 74.99 \\
F & 82.5 & 82.34& 80.33& 74.19& 74.99& 0     \\
\bottomrule
\end{tabular}
\end{center}

\textbf{Step 2: Core Distance} ($m = 2$, i.e.\ 2nd nearest neighbor):
\begin{center}
\begin{tabular}{cc}
\toprule
Point & Core Distance \\
\midrule
A & 1.41 \\
B & 1    \\
C & 1.41 \\
D & 7.07 \\
E & 7.21 \\
F & 74.99 \\
\bottomrule
\end{tabular}
\end{center}

\begin{notebox}
\textbf{Note:} If the core distance is larger than the actual distance, it means the point is sparse (itself). A sparse point cannot form a suspiciously cheap bridge to a dense group.
\end{notebox}

\textbf{Step 3: Mutual Reachability Distance (MRD)}

\[
\text{MRD}(p, q) = \max\bigl(\text{core}(p),\; \text{core}(q),\; \text{dist}(p, q)\bigr)
\]

Selected MRD values:
\begin{align*}
M(A,B) &= \max(1.41, 1, 1) = 1.41 \\
M(A,C) &= \max(1.41, 1.41, 1.41) = 1.41 \\
M(A,D) &= 8.5,\quad M(A,E) = 8.6,\quad M(A,F) = 82.5 \\
M(B,C) &= 1.41 \quad \text{(forms cycle)} \\
M(B,D) &= 7.81,\quad M(B,E) = 7.81,\quad M(B,F) = 82.34 \\
M(C,D) &= 7.07,\quad M(C,E) = 7.21,\quad M(C,F) = 80.33 \\
M(D,E) &= 7.21 \quad \text{(forms cycle)} \\
M(D,F) &= 74.99,\quad M(E,F) = 74.99
\end{align*}

\textbf{Why MRD?} To make clustering more stable, density-aware, and robust to varying densities. Direct distance can be misleading when one point is in a dense region and another is in a sparse region. MRD combines actual distance and local density information.

\textbf{Step 4: Construct MST} (keep adding lowest MRD without forming a cycle):
\[
A \xrightarrow{1.41} B,\quad A \xrightarrow{1.41} C,\quad C \xrightarrow{7.07} D,\quad C \xrightarrow{7.21} E,\quad D \xrightarrow{74.99} F
\]

\begin{center}
\begin{tabular}{ccc}
\toprule
Edge & Distance & $\lambda = 1/\text{dist}$ \\
\midrule
A--B & 1.41 & 0.71 \\
A--C & 1.41 & 0.71 \\
C--D & 7.21 & 0.14 \\
C--E & 7.07 & 0.14 \\
D--F & 74.99 & 0.01 \\
\bottomrule
\end{tabular}
\end{center}

\textbf{Step 5: Cut according to $\varepsilon$} (threshold $= 7$, cutting $D$--$F$):

Cluster: $\{A, B, C\}$; Noise: $\{D\}, \{E\}, \{F\}$ (single clusters)

\textbf{Step 6: Checking Stability of Clusters}

Density parameter $\lambda = \dfrac{1}{\text{distance}}$

Stability:
\[
S = |C| \cdot (\lambda_{\text{death}} - \lambda_{\text{birth}})
\]

\begin{itemize}
  \item Hierarchy: $\{A,B,C,D,E,F\}$ $\xrightarrow{\text{cut }74.99}$ $\{A,B,C,D,E\}$ + $\{F\}$
  \item $\{A,B,C,D,E\}$ $\xrightarrow{\text{cut }7.21}$ $\{A,B,C,E\}$ + $\{D\}$
  \item $\{A,B,C,E\}$ $\xrightarrow{\text{cut }7.07}$ $\{A,B,C\}$ + $\{E\}$
  \item $\{A,B,C\}$ $\xrightarrow{\text{cut }1.41}$ $\{A,B\}$ + $\{C\}$
  \item $\{A,B\}$ $\xrightarrow{\text{cut }1.41}$ $\{A\}$ + $\{B\}$
\end{itemize}

For $\{A,B,C\}$ cluster:
\[
\lambda_{\text{birth}} = \frac{1}{7.07} = 0.14, \quad \lambda_{\text{death}} = \frac{1}{1.41} = 0.71, \quad \text{size} = 3
\]
\[
S = 3 \times (0.71 - 0.14) = 1.71 \quad \text{(more stable)}
\]

For $\{A,B\}$:
\[
S = 2\left(\frac{1}{1.41} - \frac{1}{1.41}\right) = 0 \quad \text{(immediately splits)}
\]

For $\{A,B,C,E\}$:
\[
S = 4 \times \left(\frac{1}{7.07} - \frac{1}{7.21}\right) \approx 0.011
\]

For $\{A,B,C,D,E\}$:
\[
S = 5 \times \left(\frac{1}{7.21} - \frac{1}{74.99}\right) \approx 0.62
\]

\textbf{Result:} $\{A,B,C\}$ is the cluster; others are noise.
\end{examplebox}

\subsection{Disadvantages of H-DBSCAN}
\begin{itemize}
  \item Time consuming.
  \item Complex algorithm.
\end{itemize}

%==============================================================
\section{K-Medoids Clustering}
%==============================================================

\begin{definitionbox}[K-Medoids]
\begin{itemize}
  \item Uses \textbf{Manhattan distance}: $M(x,y) = |x_1 - x_2| + |y_1 - y_2|$.
  \item $k$ will be given.
  \item Find two medoids' values that belong to value $k_1$ and $k_2$.
\end{itemize}
\end{definitionbox}

\subsection{Example}

\begin{examplebox}
Points: $A(2,3),\; B(5,7),\; C(8,3),\; D(3,5),\; E(7,2),\; F(6,8)$,\\
$G(1,4),\; H(4,6),\; I(9,5),\; J(5,4)$. $k = 2$.

Let $A$ be medoid $k_1$ and $F$ be medoid $k_2$.

\begin{center}
\begin{tabular}{cccc}
\toprule
Point & Dist $k_1$ (A) & Dist $k_2$ (F) & Cluster \\
\midrule
A & 0 & 9  & $k_1$ \\
B & 7 & \textbf{2}  & $k_2$ \\
C & 6 & 7  & $k_1$ \\
D & \textbf{3} & 6  & $k_1$ \\
E & \textbf{6} & 7  & $k_1$ \\
F & 9 & \textbf{0}  & $k_2$ \\
G & \textbf{2} & 9  & $k_1$ \\
H & 5 & \textbf{4}  & $k_2$ \\
I & 9 & \textbf{6}  & $k_2$ \\
J & \textbf{4} & 5  & $k_1$ \\
\bottomrule
\end{tabular}
\end{center}

\textbf{Clusters:} $C_1(A,C,D,E,G,J)$ and $C_2(B,F,H,I)$

$\text{cost} = 6 + 0 + 2 + 3 + 6 + 2 + 4 \ldots = 33$ (sum minimum distances)

$\Rightarrow$ Then choose another medoid, e.g.\ $k_2 = D$, cost $= 35$ (larger, so discard; previous was good).
\end{examplebox}

\subsection{Disadvantages of K-Medoids}
\begin{itemize}
  \item Too much comparison.
  \item Time consuming.
  \item Must specify $k$.
  \item Slower convergence.
\end{itemize}

%==============================================================
\section{Hierarchical Clustering (Agglomerative)}
%==============================================================

\subsection{Example}

\begin{examplebox}
Points: $P_1(0.40,0.53),\; P_2(0.22,0.38),\; P_3(0.35,0.32),\; P_4(0.26,0.19),\; P_5(0.08,0.41),\; P_6(0.45,0.30)$

\textbf{Distance Matrix:}
\begin{center}
\begin{tabular}{ccccccc}
\toprule
 & $P_1$ & $P_2$ & $P_3$ & $P_4$ & $P_5$ & $P_6$ \\
\midrule
$P_1$ & 0    &      &       &       &       &       \\
$P_2$ & 0.23 & 0    &       &       &       &       \\
$P_3$ & 0.22 & 0.15 & 0     &       &       &       \\
$P_4$ & 0.37 & 0.20 & 0.15  & 0     &       &       \\
$P_5$ & 0.34 & 0.14 & 0.28  & 0.29  & 0     &       \\
$P_6$ & 0.23 & 0.25 & \textbf{0.11} & 0.22 & 0.39 & 0     \\
\bottomrule
\end{tabular}
\end{center}

\textbf{Iteration 1:} Minimum distance $= 0.11$ between $P_3$ and $P_6$. Merge $\to P_{3,6}$.

\textbf{Iteration 2:} Update matrix; minimum $= 0.14$ between $P_2$ and $P_5$. Merge $\to P_{2,5}$.

\textbf{Iteration 3:} Minimum $= 0.15$ between $P_{3,6}$ and $P_4$. Merge $\to P_{3,4,6}$.

\textbf{Iteration 4:} Minimum $= 0.22$ between $P_1$ and $P_{2,3,4,5,6}$. Merge all.

\textbf{Dendrogram:} Shows nested groupings: $\{2,5\},\;\{3,6\} \to \{3,4,6\} \to$ all together.

The diagram shows: Points 2 and 5 merge first, 3 and 6 merge, then 3,6,4 merge, then all points join at the top.
\end{examplebox}

%==============================================================
\section{Decision Trees (ID3 Algorithm)}
%==============================================================

\subsection{Entropy}

\[
\text{Entropy}(S) = -\sum_i p_i \log_2(p_i)
\]

\subsection{Information Gain}

\[
\text{Gain}(S, a) = \text{Entropy}(S) - \sum_v \frac{|S_v|}{|S|}\,\text{Entropy}(S_v)
\]

\subsection{Example}

\begin{examplebox}
Dataset ($n=10$): Attributes $a_1$ (T/F), $a_2$ (H/C), $a_3$ (H/N), Classification (Yes/No).

\begin{center}
\begin{tabular}{ccccc}
\toprule
ID & $a_1$ & $a_2$ & $a_3$ & Class \\
\midrule
1 & T & H & H & No \\
2 & T & H & H & No \\
3 & F & H & H & Yes \\
4 & F & C & N & Yes \\
5 & F & C & N & Yes \\
6 & T & C & H & No \\
7 & T & H & H & No \\
8 & T & H & N & Yes \\
9 & F & C & N & Yes \\
10 & F & C & H & Yes \\
\bottomrule
\end{tabular}
\end{center}

$S = [6+, 4-]$:
\[
\text{En}(S) = -\frac{6}{10}\log_2\!\left(\frac{6}{10}\right) - \frac{4}{10}\log_2\!\left(\frac{4}{10}\right) = 0.9709
\]

$S_\text{True} = [1+, 5-]$:
\[
E(S_\text{True}) = -\frac{1}{5}\log_2\!\left(\frac{1}{5}\right) - \frac{4}{5}\log_2\!\left(\frac{4}{5}\right) = 0.7219
\]

$S_\text{False} = [5+, 0-]$:
\[
E(S_\text{False}) = 0.0
\]

\[
\text{Gain}(S, a_1) = 0.9709 - \frac{5}{10}(0.7219) - \frac{5}{10}(0.0) = 0.6099
\]
\end{examplebox}

%==============================================================
\section{Logistic Regression}
%==============================================================

\subsection{Background: Linear Regression}

\[
\hat{y} = m_1 x_1 + m_2 x_2 + m_3 x_3 + \cdots
\]

\subsection{Sigmoid Function}

The sigmoid function comes from the idea of converting any real number into a value between 0 and 1, so that it can represent probability.

In logistic regression: $P(y=1) = ?$, $0 \leq P \leq 1$.

But the linear equation $z = w_1 x_1 + w_2 x_2 + b$ can produce $-\infty$ to $+\infty$.

\[
\text{odds} = \frac{p}{1-p} \quad \text{where } p = \text{probability of class 1}
\]

If $p = 0.8$: $\text{odds} = \dfrac{0.8}{0.2} = 4$ (it means success is 4 times more likely).

\[
\text{logit} = \log\!\left(\frac{p}{1-p}\right) \quad \text{(can range from } -\infty \text{ to } +\infty\text{)}
\]

\textbf{Derivation of sigmoid:}
\[
\log\!\left(\frac{p}{1-p}\right) = z \quad \Rightarrow \quad \frac{p}{1-p} = e^z \quad \Rightarrow \quad p = e^z - e^z p
\]
\[
\Rightarrow\quad p(1 + e^z) = e^z \quad \Rightarrow \quad
\boxed{p(y) = \frac{1}{1+e^{-z}}} \quad \text{(sigmoid function)}
\]

\begin{itemize}
  \item It transforms probabilities into an unbounded scale, making them suitable for linear modeling.
  \item The sigmoid function makes the output of linear regression bounded.
  \item The value after passing through the activation function is \textbf{not} the actual probability.
\end{itemize}

%==============================================================
\section{Loss Function}
%==============================================================

\begin{definitionbox}[Loss Function]
\begin{itemize}
  \item Measures how wrong the model's prediction is.
  \item Smaller loss means better prediction.
  \item ML algorithm tries to \textbf{minimize the loss}.
\end{itemize}
\end{definitionbox}

Let $y$ = actual/expected output, $\hat{y}$ = actual output. Then $y - \hat{y} = 0$ means no error.

\subsection{Binary Cross-Entropy Loss}

\[
\boxed{\mathcal{L} = -\left[y\ln(\hat{y}) + (1-y)\ln(1-\hat{y})\right]}
\]

\textbf{Cases:}

\begin{enumerate}
  \item \textbf{$y = \hat{y} = 0$:}
        \[
        \mathcal{L} = -\bigl[0 + (1-0)\ln(1-0)\bigr] = -\ln(1) = 0 \quad \text{(no error)}
        \]

  \item \textbf{$y = 0,\; \hat{y} = 1$} (worst possible prediction):
        \[
        \mathcal{L} = -\bigl[0 + (1-0)\ln(1-1)\bigr] = -\ln(0) = \infty \quad \text{(extremely large)}
        \]

  \item \textbf{$y = 1,\; \hat{y} = 0$}:
        \[
        \mathcal{L} = -\bigl[1 \cdot \ln(0) + 0\cdot \ln(1)\bigr] = -\ln(0) = \infty \quad \text{(extremely large)}
        \]

  \item \textbf{$y = \hat{y} = 1$:}
        \[
        \mathcal{L} = -\bigl[1\cdot \ln(1) + 0\cdot\ln(0)\bigr] = 0 \quad \text{(correct prediction)}
        \]
\end{enumerate}

%==============================================================
\section{Neural Networks}
%==============================================================

\begin{definitionbox}[Neural Networks]
\begin{itemize}
  \item A machine learning model.
  \item It learns patterns from data using interconnected units called \textbf{neurons}.
  \item Usually has input, hidden, and output layers.
  \item Used for both regression and classification.
  \item Neuron: $x_i \xrightarrow{w} \bigcirc \xrightarrow{} y$, where $y = wx + b$.
  \item Neurons work linearly, but that is not effective. That's why \textbf{non-linearity is introduced by activation functions}.
  \item Common activations: Sigmoid $\to \dfrac{1}{1+e^{-z}}$,\quad ReLU $\to \max(0,x)$,\quad $\tanh \to \tanh(z)$.
\end{itemize}
\end{definitionbox}

%==============================================================
\section{Autodiff (Automatic Differentiation)}
%==============================================================

\subsection{Normal way vs.\ Autodiff}

Let $a=4,\; b=3,\; c=a+b,\; d=a \times c$.

\textbf{Traditional way:}
\[
d = a \times c \quad \Rightarrow \quad \frac{\partial d}{\partial a} = \frac{\partial a}{\partial a}\cdot c + \frac{\partial c}{\partial a}\cdot a = c + a \cdot \frac{\partial c}{\partial a}
\]
\[
= (a+b) + a \cdot \frac{\partial(a+b)}{\partial a} = a+b + a\cdot\left(\frac{\partial a}{\partial a} + \frac{\partial b}{\partial a}\right) = a+b+a(1+0)
\]
\[
= a + b + a = 4+3+4 = 11 \text{\checkmark}
\]

\textbf{Autodiff} (local derivatives):
Local derivatives through the computation graph:
\[
\frac{\partial d}{\partial a} = \underbrace{\frac{\partial d}{\partial c}\cdot\frac{\partial c}{\partial a}}_{\text{path via }c} + \underbrace{\frac{\partial d}{\partial a}}_{\text{direct}}
= a \cdot 1 + c = a + (a+b) = a + b + a = 11
\]

\subsection{Comparison: Local vs.\ Partial Derivative}

\begin{center}
\begin{tabular}{p{6cm}p{6cm}}
\toprule
\textbf{Local Derivative} & \textbf{Partial Derivative} \\
\midrule
Derivative at one operation & Derivative in multivariable function \\
Immediate relationship & One variable fixed, others change \\
E.g.\ $z = x^2 \Rightarrow \dfrac{dz}{dx} = 2x$ & E.g.\ $f(x,y) = x^2 + 3y \Rightarrow \dfrac{\partial f}{\partial x} = 2x,\; \dfrac{\partial f}{\partial y} = 3$ \\
Only includes $z$ and $x$ & Multiple variables \\
\bottomrule
\end{tabular}
\end{center}

\subsection{Use of Autodiff}
\begin{itemize}
  \item Computes $\dfrac{\partial \mathcal{L}}{\partial w}$ for all weights during \textbf{backpropagation}.
  \item Backpropagation uses the \textbf{reverse mode} of autodiff.
\end{itemize}

%==============================================================
\section{Autoencoder}
%==============================================================

\begin{definitionbox}[Autoencoder]
A type of neural network architecture used in DL for \textbf{unsupervised learning} and \textbf{dimensionality reduction} tasks.
\begin{itemize}
  \item Compresses data and makes encoded data.
  \item Keeps the important features.
  \item Decoding means reconstruction.
\end{itemize}
\end{definitionbox}

\subsection{Structure}

Encoder: $512 \to 256 \to 128 \to 64 \to \mathbf{49}$ (bottleneck) $\to 64 \to 128 \to 256 \to 512$: Decoder

\begin{itemize}
  \item Auto-encoder uses a bottleneck layer which forces the network to learn compact dimensionality representation. It uses the $64 \times (h \times w)$ and other dimension features.
  \item Compressed: $1000 \to 500$.
  \item Auto-encoder both the input take. Input depth is smaller.
\end{itemize}

%==============================================================
\section{NN for Dimensionality Reduction}
%==============================================================

NN mainly uses for dimensionality reduction. For both the input, auto-encoder takes the input depth.

The encoder part reduces dimensions; the decoder part (reconstruction) tries to recreate the original input.

%==============================================================
\section{NN for Outlier Detection}
%==============================================================

Here, \textbf{auto-encoder is used again}.

\begin{enumerate}
  \item Train the autoencoder using \textbf{normal data}. The network becomes good at reconstructing the normal patterns.
  \item If it gives bad reconstruction, then \textbf{unusual data were used}.
  \item Reconstruction error (loss):
        \[
        \mathcal{L} = \|x - \hat{x}\|^2
        \]
        where $x$ = original input, $\hat{x}$ = reconstructed output.
  \item If error is small $\Rightarrow$ no anomaly; if error is larger $\Rightarrow$ outlier.
\end{enumerate}

\subsection{Example}

\begin{examplebox}
$x = [1, 2, 3]$, $\hat{x} = [1.1, 2.0, 2.9]$

\[
\mathcal{L} = (0.1)^2 + (0)^2 + (0.1)^2 = 0.02 \quad \text{(small)} \Rightarrow \text{not an outlier}
\]
\end{examplebox}

%==============================================================
\section{Why Normal ANN is Not Suitable for Image Processing}
%==============================================================

\textbf{Reason 1: Too Many Parameters}

A small $224 \times 224$ RGB image has $3 \times 224^2 = 150{,}528$ input values. If the first hidden layer has just 1000 neurons, that's $\sim 150$ million parameters --- just for one layer. This leads to:
\begin{itemize}
  \item Enormous memory requirements.
  \item Slow training.
\end{itemize}

\textbf{Reason 2: No Spatial Awareness}

Plain NN treats every pixel as an independent input. They have no concept of spatial relationship. They lose all 2D structure by flattening the image into a 1D vector.

\textbf{Reason 3: Not Translation Invariant}

It has no relation for the same concept at every possible location. One pen at left corner, another pen at bottom right --- NN treats them completely differently.

\textbf{That's why NN is not suitable $\Rightarrow$ use CNN.}

%==============================================================
\section{CNN (Convolutional Neural Network)}
%==============================================================

\begin{definitionbox}[CNN]
\begin{itemize}
  \item Convolutional Neural Network.
  \item Designed to understand \textbf{edges, textures, shapes, objects} from images.
  \item Uses convolution, filters, pooling.
  \item Has many layers.
\end{itemize}
\end{definitionbox}

\subsection{Layers}

\textbf{Convolution Layer:} Uses small filters/kernels, e.g.:
\[
\begin{pmatrix} 1 & 0 & -1 \\ 1 & 0 & -1 \\ 1 & 0 & -1 \end{pmatrix} \quad \text{(detects vertical edges)}
\]

\textbf{Pooling Layer:} Reduces image size. Most common: max pooling, average pooling, etc.

\textbf{Fully Connected Layer:} Final classification stage.

\subsection{Working Flow}

\[
\text{Image} \to \text{Convolution} \xrightarrow{\text{ReLU}} \text{Pooling} \to \text{repeat} \to \text{Fully Connected} \to \text{Prediction}
\]

%==============================================================
\section{Excitation Channels (SE Block)}
%==============================================================

Not every channel has the same importance. So we assign weights. If one channel doesn't contain any info, its weight will be 0. Per-channel attention will be calculated by MPN (max pooling network).

\begin{definitionbox}[Squeeze-and-Excitation (SE) Block]
\textbf{Channel attention} means the network learns which feature channels are more important. One famous method is the \textbf{Squeeze-and-Excitation (SE) Block}.

After convolution, CNN produces many feature maps/channels (e.g.\ edge detector channel, texture channel, color channel). But not every channel is equally important.

SE blocks learn which channels should be \textbf{emphasized} and which should be \textbf{suppressed}.
\end{definitionbox}

\subsection{Steps}

\textbf{Step 1 -- Squeeze:} Compress each feature map into one number using \textbf{Global Average Pooling}:
\[
z_c = \frac{1}{H \times W} \sum_{i=1}^{H}\sum_{j=1}^{W} x_c(i,j)
\]
This captures the importance of a channel.

\textbf{Step 2 -- Excitation:} The compressed value passes through a small neural network (sigmoid functions are used).

$\Rightarrow$ Finally multiply channel by attention weight.

\subsection{Workflow}

\[
\text{Feature map} \xrightarrow{\text{Squeeze (Global Avg Pooling)}} \xrightarrow{\text{Excitation (ReLU + Sigmoid)}} \text{Attention weights}
\]
\[
\xrightarrow{\text{Channel reweighting}} \text{Refined feature map}
\]

%==============================================================
\section{Depthwise Separable Convolution}
%==============================================================

\begin{definitionbox}[Depthwise Separable Convolution]
\begin{itemize}
  \item Reduces the number of floating point operations and learnable parameters.
  \item Splits convolution into 2 parts:
        \begin{enumerate}
          \item Depthwise convolution
          \item Pointwise convolution
        \end{enumerate}
\end{itemize}
\end{definitionbox}

\subsection{Normal Convolution vs.\ DSC}

\textbf{Normal Convolution:}\\
Filter: $5\times5\times64$, Input: $32\times32\times64$, Output: $32\times32\times128$
\[
\text{\# learnable params} = 5\times5\times64\times128 = 204{,}800 \approx 204\text{k}
\]
\[
\text{\# multiplications} = 5\times5\times64\times32\times32\times128 \approx 209\text{ Million}
\]

\textbf{Depthwise Separable Convolution (DSC):}\\
Filter: $5\times5\times1$ (per channel), Input: $32\times32\times64$
\[
\text{\# learnable params} = (5\times5\times1\times64) + (1\times1\times64\times128) = 1600 + 8192 = 9792
\]
\[
\text{\# multiplications} = (5\times5\times1\times32\times32\times64) + (1\times1\times64\times32\times3\times64)
\approx 10\text{M}
\]

$\Rightarrow$ Thus we're able to \textbf{significantly reduce} the number of learnable parameters and multiplications. So less memory and computation is required.

\begin{center}
\begin{tabular}{lccc}
\toprule
Feature & Normal Conv & Depthwise & Dilated \\
\midrule
Goal & Feature extraction & Reduce computation & Increase receptive field \\
Parameters & High & Low & Similar to normal \\
Speed & Slower & Faster & Moderate \\
Used in & Standard CNN & Mobile CNN & Segmentation context \\
\bottomrule
\end{tabular}
\end{center}

%==============================================================
\section{Dilated Convolution}
%==============================================================

\begin{definitionbox}[Dilated Convolution]
A technique that expands the kernel (input) by inserting holes between its consecutive elements. It is the same as convolution but it involves pixel skipping, so as to cover a larger area of the input.

\[
\Rightarrow \text{Without increasing the parameters, it enables the network to have a larger \textbf{receptive field}.}
\]
\end{definitionbox}

Dilation factor $l=1$ (normal convolution) vs.\ $l=2$ (dilated convolution, skips every other pixel).

%==============================================================
\section{Outlier Detection Methods}
%==============================================================

\begin{enumerate}
  \item \textbf{Statistical method}: Z-score ($|z| > 3$ then outlier)
  \item \textbf{Distance-based method}: KNN
  \item \textbf{Density-based method}: DBSCAN
  \item \textbf{Neural network method}: Autoencoder
\end{enumerate}

%==============================================================
\section{Problem with Euclidean Distance}
%==============================================================

\begin{notebox}
Euclidean distance does not consider \textbf{correlation, scale, covariance}.
\end{notebox}

\subsection{Example}

\begin{examplebox}
Height-weight data:
\begin{center}
\begin{tabular}{cc}
\toprule
Height (cm) & Weight (kg) \\
\midrule
150 & 45 \\
160 & 55 \\
170 & 65 \\
180 & 75 \\
\bottomrule
\end{tabular}
\end{center}

Strong correlation: taller people have heavier weight.

Now if Height, Weight $= (180, 45)$: individually 180 and 45 are normal, but \textbf{together} abnormal (heavyweight is low for that height).

$\Rightarrow$ Euclidean fails to predict this.
\end{examplebox}

%==============================================================
\section{Mahalanobis Distance}
%==============================================================

\begin{definitionbox}[Mahalanobis Distance]
It considers \textbf{covariance, correlation between features, multivariant relationships}.

\[
\boxed{D_M = \sqrt{(x-\mu)^T \,\Sigma^{-1}\, (x-\mu)}}
\]

where $\Sigma^{-1}$ is the inverse covariance matrix.

\textbf{Larger $D_M$ means outlier.}
\end{definitionbox}

\subsection{Decision Rule}
\[
D_M^2 > \chi^2_{\text{threshold}} \Rightarrow \text{anomaly}
\]
\[
D_M^2 < \chi^2_{\text{threshold}} \Rightarrow \text{normal}
\]

\subsection{Example}

\begin{examplebox}
Student height-weight data:
\begin{center}
\begin{tabular}{ccc}
\toprule
Student & Height (cm) & Weight (kg) \\
\midrule
A & 150 & 50 \\
B & 160 & 60 \\
C & 170 & 70 \\
D & 180 & 80 \\
E & 180 & 50 \\
\bottomrule
\end{tabular}
\end{center}
Check if $E(180, 50)$ is an anomaly or not.

\textbf{Solution:}
\[
\mu_h = 168, \quad \mu_w = 62, \quad \bm{\mu} = (168, 62)
\]
\[
\Sigma = \begin{pmatrix} 170 & 110 \\ 110 & 170 \end{pmatrix}, \quad
\Sigma^{-1} = \begin{pmatrix} 0.0107 & -0.0069 \\ -0.0069 & 0.0107 \end{pmatrix}
\]

For $E$: $x - \mu = \begin{pmatrix} 12 \\ -12 \end{pmatrix}$

\[
D_M^2 = \begin{pmatrix} 12 & -12 \end{pmatrix}
\begin{pmatrix} 0.0107 & -0.0069 \\ -0.0069 & 0.0107 \end{pmatrix}
\begin{pmatrix} 12 \\ -12 \end{pmatrix} = 5.30
\]
\[
D_M = \sqrt{5.30} \approx 2.30
\]

Degrees of freedom $= 2$.

$\chi^2_{0.95,\;2} = 5.99$

Since $D_M^2 = 5.30 < 5.99$, point $E$ is \textbf{normal}.
\end{examplebox}

\end{document}